\section{Periods and intersections of distinguished subgroups}\label{periods-and-intersections-of-distinguished-subgroups}

Let \(w_1,\ldots,w_d\in U=E\) be translation words with linearly independent
vectors \(v_1,\ldots,v_d\). Write
\(\Lambda=\sum_i\mathbb Z v_i\leq L\) and suppose it is \(Q\)-invariant.
For a positive integer \(m\), adjoining \(w_i^m=1\) to the complete
presentation of \(U\) gives exactly

\[ E_m=E/(m\Lambda).\tag{3}\]

Indeed, the normal closure of the powers is \(m\Lambda\): conjugation by
an affine element acts on translations by its linear part, and this
lattice is already invariant. Hence

\[ |E_m|=|Q|[L:\Lambda]m^d.\tag{4}\]

This proves equality for the presented group, not merely the order of a
chosen quotient. If \(Q\) is solvable, every \(E_m\) is solvable.

Suppose now that the marking has rank three with distinct generators,
its endpoints commute, both adjacent products have order at least three, and
its two adjacent subgroups \(A=E_{01}\) and \(B=E_{12}\) are finite and
inject into \(Q\). For each \(q\in\overline A\cap\overline B\), let
\((q,t_A(q))\) and \((q,t_B(q))\) be their unique affine lifts. Put

\[ \Delta(q)=t_A(q)-t_B(q).\tag{5} \]

\textbf{Theorem AL.3 (period-lattice intersection criterion).} The images of \(A\)
and \(B\) in \(E_m\) meet precisely in the elements indexed by

\[ \{q\in\overline A\cap\overline B:\Delta(q)\in m\Lambda\}.\tag{6}\]

In particular, if \(A\cap B=\langle r_1\rangle\) and every other value
\(\Delta(q)\) lies outside \(\Lambda\), then every \(E_m\), \(m\geq1\), has
the full rank-three intersection property and the same adjacent orders
as \(E\). Its order is (4).

\emph{Proof.} Equality of two cosets modulo translations first forces equality
of their linear parts. For equal linear part, their affine quotient is
translation by the difference (5). This proves (6). Each adjacent
subgroup remains injective in the quotient because its linear action is
faithful, so its order and the product order of its two reflections are
unchanged. The assumed exclusions apply to \(m\Lambda\subseteq\Lambda\).
Thus the adjacent intersection is exactly the common involution subgroup,
and the rank-three test proves all subset intersections. \(\square\)

The proof separates two possibilities often hidden by equal group-order
counts. The point-group intersection may already have order two, as in
the \(\{6,4\}\) application. Or it may be larger, with a short list of
nonzero translation differences, as in the \(\{6,6\}\) application. Those
differences must be excluded from the period lattice explicitly.

\subsection*{One coarse quotient controls all smaller periods}
The same argument has a group-theoretic formulation which is useful independently of coordinates. Let $A,B,L$ be subgroups of a group $E$, with
\[
 L\cap A=L\cap B=1.
\]
Let $N\unlhd E$ lie in $L$. If the images of $A,B$ in $E/N$ intersect exactly in the image of $A\cap B$, then the same statement holds in $E/M$ for every normal subgroup $M\le N$.

Indeed, suppose $a\in A$, $b\in B$ have the same image modulo $M$. They have the same image modulo $N$, so there is $c\in A\cap B$ with $aN=cN=bN$. Then $ac^{-1}\in N\cap A=1$ and $bc^{-1}\in N\cap B=1$, proving $a=b=c$. The reverse inclusion is immediate. In the affine setting $L$ is the translation kernel, $N=\Lambda$ and $M=m\Lambda$. Point-faithfulness supplies precisely the two trivial intersections. This statement, for an arbitrary indexed family of smaller normal periods, is verified in \texttt{AffinePeriodIntersections.lean}. The finite coordinate checks which identify $L,\Lambda$ and the initial overlap remain separate mathematical obligations.
